\phantomsection
\addcontentsline{toc}{chapter}{\bf ABSTRACT}

\begin{center}
    \large{\bf ABSTRACT}
\end{center}

The growing demand for psychological support has increased interest in using large language models (LLMs) as accessible counseling agents. However, therapeutic dialogue requires more than generating a locally appropriate response. In counseling, especially when clients show resistance or ambivalence, the timing of therapeutic strategies is critical. Motivational Interviewing (MI) is useful for supporting readiness, autonomy, and exploration of ambivalence, while Cognitive Behavioral Therapy (CBT) provides structured techniques for identifying and modifying unhelpful thoughts and behaviors. Introducing CBT too early may increase resistance, while remaining only in MI may slow progress when the client is ready for structured intervention.

Existing LLM-based counseling systems have explored therapeutic strategy selection, memory-based adaptation, and MI-CBT dialogue support. However, many approaches remain limited in their ability to anticipate the immediate downstream effects of a therapeutic response before it is selected. As a result, a system may choose a response that appears suitable in the current turn but does not support the client’s next therapeutic state, readiness, or engagement.

This thesis proposes AIM-RGL, an agentic MI-CBT framework with reward-guided lookahead search for managing client resistance in simulated therapeutic conversations. The framework uses multiple specialist therapist agents, including open-ended questioning, reflection, affirmation, summarization, and CBT agents, to generate candidate responses. A client simulator then produces possible future client reactions based on cognitive conceptualization and stage-of-change information. A reward model evaluates candidate responses according to therapeutic dimensions such as alliance, empathy, collaboration, CBT skill, MI quality, and client readiness. Instead of selecting only the best immediate response, AIM-RGL performs short-horizon lookahead search to choose the response expected to lead to a better near-future therapeutic state.

The study evaluates whether reward-guided lookahead improves MI-CBT response quality compared with baseline approaches. It further examines how different lookahead depths, reward model backbones, and policy model backbones affect therapeutic response generation. Evaluation focuses on simulated therapist-client dialogues using metrics related to therapeutic alliance, MI fidelity, CBT competence, change talk, and sustain talk. This thesis does not claim to measure real clinical treatment outcomes or replace professional therapists; rather, it investigates how future-aware response selection can improve the design of LLM-based counseling support systems. Overall, AIM-RGL contributes a structured framework for proactive, context-sensitive, and therapeutically grounded response generation in AI-assisted mental health dialogue.
